% --- IMPORTS ---
\documentclass[leqno]{article}
\usepackage{graphicx}
\usepackage{dsfont}
\usepackage{mathtools}
\usepackage{amssymb}
\usepackage{amsmath}
\usepackage{amsthm}
\usepackage{float}
\usepackage{ragged2e}
\usepackage{tensor}
\usepackage{fancyhdr}
\usepackage{empheq}
\usepackage[most]{tcolorbox}
\usepackage{enumitem}
\usepackage{dirtytalk}
\usepackage{marginnote}
\usepackage{microtype}
\usepackage[
  a4paper,
  left=0.8in,
  right=1in,
  top=1in,
  bottom=0.5in,
  headheight=14pt,
  headsep=0.4in
]{geometry}
\setlength{\parindent}{0cm}
% --- TikZ Packages ---
\usepackage{pgfplots}
\pgfplotsset{compat=1.18}
\usepgfplotslibrary{fillbetween, polar}
\usetikzlibrary{calc, positioning}
\usetikzlibrary{angles, quotes}
\usetikzlibrary{arrows.meta}
\usetikzlibrary{decorations.pathreplacing, decorations.markings}
\usetikzlibrary{patterns}
\usetikzlibrary{intersections}
% --- URL Packages ---
\usepackage[colorlinks=true,linkcolor=red,citecolor=magenta,urlcolor=black]{hyperref}%
\urlstyle{same}
% --- END OF IMPORTS ---

% --- CUSTOM COMMANDS ---
\RenewDocumentCommand{\marks}{o m}{%
  \IfValueTF{#1}%
    {\marginnote{\raggedright\textbf{[#2]}}[#1]}%
    {\marginnote{\raggedright\textbf{[#2]}}}%
}
\newcommand{\vspacerule}[2]{%
  \par\vspace*{#1}%
  \noindent\hrulefill\par
  \vspace*{#2}%
}
\definecolor{scarlet}{RGB}{205, 40, 75}
\newcommand{\questionitem}{%
  \item\phantomsection
  \addcontentsline{toc}{section}{Question \number\value{enumi}}%
}
\renewcommand{\contentsname}{Questions}
% --- END OF CUSTOM COMMANDS ---

% --- HEADER CONFIG ---
\pagestyle{fancy}
\fancyhead{}
\fancyfoot{}
\renewcommand{\headrulewidth}{0.9pt}
\lhead{\textit{Solutions} - Invariant Points and Lines}
\chead{}
\rhead{\href{https://danieljsmith.org}{danieljsmith.org}}
% --- END OF HEADER CONFIG ---

\begin{document}
\vspace*{20pt}
\tableofcontents
\newpage
\begin{enumerate}[
  label=\textbf{\arabic*.},
  leftmargin=*,
  topsep=0pt,
  partopsep=0pt,
  parsep=0pt
]

% ---- Question 1 ----
\questionitem
% [Source: _QBT___Solns__Invariant_Points_and_Lines, Q1]

The matrix $M$ represents the transformation $T$ and is given by
\[
M = \begin{pmatrix}
p + 3 & -2 \\[6pt]
 p + 2 & 1 - p
\end{pmatrix}
\] \\
\begin{enumerate}[label=\textbf{(\alph*)}]
  \item In the case when $p = -1$ show that the image of the point $(3, 4)$ under $T$ is the point $(-2, 11)$. \marks{2} \\
  \item In the case when $p = 2$ find the gradient of the line of invariant points under $T$. \marks{3} \\
\end{enumerate}

\vspace*{10pt}

\begin{tcolorbox}[colback=white, colframe=scarlet, title={\textbf{Solution}}, breakable, grow to left by=6mm, grow to right by=10mm]
\begin{enumerate}[label=\textbf{(\alph*)}]
  \item When $p=-1$,
\[
M=\begin{pmatrix}
p+3 & -2\\
p+2 & 1-p
\end{pmatrix}
=
\begin{pmatrix}
2 & -2\\
1 & 2
\end{pmatrix}
\]

The image of the point $(3,4)$ is found by multiplying by the matrix:
\[
\begin{pmatrix}
2 & -2\\
1 & 2
\end{pmatrix}
\begin{pmatrix}
3\\
4
\end{pmatrix}
=
\begin{pmatrix}
2(3)-2(4)\\
1(3)+2(4)
\end{pmatrix}
=
\begin{pmatrix}
6-8\\
3+8
\end{pmatrix}
=
\begin{pmatrix}
-2\\
11
\end{pmatrix}
\]

So the image of $(3,4)$ under $T$ is $(-2,11)$, as required.

  \item When $p=2$,
\[
M=\begin{pmatrix}
p+3 & -2\\
p+2 & 1-p
\end{pmatrix}
=
\begin{pmatrix}
5 & -2\\
4 & -1
\end{pmatrix}
\]

Let an invariant point be $(x,y)$. Then
\[
\begin{pmatrix}
5 & -2\\
4 & -1
\end{pmatrix}
\begin{pmatrix}
x\\
y
\end{pmatrix}
=
\begin{pmatrix}
x\\
y
\end{pmatrix}
\]

This gives
\begin{align*}
5x-2y&=x\\
4x-y&=y
\end{align*}

Rearranging,
\begin{align*}
4x-2y&=0\\
4x-2y&=0
\end{align*}

So the invariant points lie on the line
\[
4x-2y=0
\]

Hence
\[
y=2x
\]

Therefore the gradient of the line of invariant points is $2$.
\end{enumerate}
\end{tcolorbox}


\newpage

% ---- Question 2 ----
\questionitem
% [Source: _QBT___Solns__Invariant_Points_and_Lines, Q2]

A linear transformation of the plane is represented by
\[
\begin{pmatrix}
1 & 4 \\
2 & 3
\end{pmatrix}
\]
Determine the invariant line with negative gradient. \marks{4}

\vspace*{10pt}

\begin{tcolorbox}[colback=white, colframe=scarlet, title={\textbf{Solution}}, breakable, grow to left by=6mm, grow to right by=10mm]
Let the invariant line have equation
\[
y=mx+c
\]

Take a general point on the line, \((x,mx+c)\). Under the transformation
\[
\begin{pmatrix}x\\y\end{pmatrix}
\mapsto
\begin{pmatrix}
1 & 4\\
2 & 3
\end{pmatrix}
\begin{pmatrix}x\\y\end{pmatrix}
\]
this point maps to
\[
\begin{pmatrix}
1 & 4\\
2 & 3
\end{pmatrix}
\begin{pmatrix}
x\\
mx+c
\end{pmatrix}
=
\begin{pmatrix}
x+4(mx+c)\\
2x+3(mx+c)
\end{pmatrix}
=
\begin{pmatrix}
(1+4m)x+4c\\
(2+3m)x+3c
\end{pmatrix}
\]

For the line to be invariant, the image point must also lie on \(y=mx+c\). So
\[
(2+3m)x+3c = m\bigl((1+4m)x+4c\bigr)+c
\]

Expanding,
\[
(2+3m)x+3c=(m+4m^2)x+(4m+1)c
\]

Comparing coefficients of \(x\) and the constant terms,
\begin{align*}
2+3m&=m+4m^2\\
3c&=(4m+1)c
\end{align*}

From the first equation,
\begin{align*}
4m^2-2m-2&=0\\
2m^2-m-1&=0\\
(2m+1)(m-1)&=0
\end{align*}

So
\[
m=1 \quad \text{or} \quad m=-\frac12
\]

From the second equation,
\[
(2-4m)c=0
\]

For both values of \(m\), this gives \(c=0\).

So the invariant lines are
\[
y=x
\quad \text{and} \quad
y=-\frac12x
\]

Therefore, the invariant line with negative gradient is
\[
y=-\frac12x
\]

Equivalently, \(x+2y=0\).
\end{tcolorbox}


\newpage

% ---- Question 3 ----
\questionitem
% [Source: _QBT___Solns__Invariant_Points_and_Lines, Q3]

The matrix $A$ is given by
\[
A = \frac{1}{5}\begin{pmatrix}-3 & -4 \\ -4 & 3\end{pmatrix}
\] \\
You are given that $A$ represents the transformation $T$, where $T$ is a reflection in a straight line through the origin, $O$. \\
\begin{enumerate}[label=\textbf{(\alph*)}]
  \item Explain why $T$ must have a line of invariant points. State the geometric significance of this line. \marks{2} \\
  \item By considering the line of invariant points for $T$, determine the equation of the mirror line.

Give your answer in the form $y = mx + c$. \marks{4} \\
  \item The coordinates of the point $P$ are $(4, 7)$. \\
  By considering the image of $P$ under the transformation $T$, or otherwise, determine the coordinates of the point on the mirror line which is closest to $P$. \marks{3} \\
  \item The line with equation $2x + by = 6$ is an invariant line for $T$. \\
  Determine the value of $b$. \marks{2} \\
\end{enumerate}

\vspace*{10pt}

\begin{tcolorbox}[colback=white, colframe=scarlet, title={\textbf{Solution}}, breakable, grow to left by=6mm, grow to right by=10mm]
\begin{enumerate}[label=\textbf{(\alph*)}]
  \item A reflection leaves every point on its mirror line unchanged, so $T$ must have a line of invariant points.

  Geometrically, this line is the mirror line (the axis of reflection).

  \item Let an invariant point be $(x,y)$. Then
  \[
  A\begin{pmatrix}x\\y\end{pmatrix}=\begin{pmatrix}x\\y\end{pmatrix}
  \]

  So
  \[
  \frac{1}{5}
  \begin{pmatrix}
  -3 & -4\\
  -4 & 3
  \end{pmatrix}
  \begin{pmatrix}x\\y\end{pmatrix}
  =
  \begin{pmatrix}x\\y\end{pmatrix}
  \]

  This gives
  \[
  \frac{1}{5}
  \begin{pmatrix}
  -3x-4y\\
  -4x+3y
  \end{pmatrix}
  =
  \begin{pmatrix}x\\y\end{pmatrix}
  \]

  Hence
  \begin{align*}
  -3x-4y&=5x\\
  -4x+3y&=5y
  \end{align*}

  So
  \begin{align*}
  -8x-4y&=0\\
  -4x-2y&=0
  \end{align*}

  These are the same equation, giving
  \[
  2x+y=0
  \]

  Therefore the mirror line is
  \[
  y=-2x
  \]

  So the equation of the mirror line is $y=-2x$.

  \item First find the image of $P=(4,7)$ under $T$:
  \[
  A\begin{pmatrix}4\\7\end{pmatrix}
  =
  \frac{1}{5}
  \begin{pmatrix}
  -3 & -4\\
  -4 & 3
  \end{pmatrix}
  \begin{pmatrix}4\\7\end{pmatrix}
  =
  \frac{1}{5}
  \begin{pmatrix}
  -12-28\\
  -16+21
  \end{pmatrix}
  =
  \frac{1}{5}
  \begin{pmatrix}
  -40\\
  5
  \end{pmatrix}
  =
  \begin{pmatrix}
  -8\\
  1
  \end{pmatrix}
  \]

  So the image is $P'=(-8,1)$.

  For a reflection, the mirror line is the perpendicular bisector of $PP'$. So the closest point on the mirror line to $P$ is the midpoint of $PP'$:
  \[
  \left(\frac{4+(-8)}{2},\frac{7+1}{2}\right)
  =
  \left(\frac{-4}{2},\frac{8}{2}\right)
  =
  (-2,4)
  \]

  Therefore the point on the mirror line closest to $P$ is $(-2,4)$.

  \item The mirror line is $y=-2x$, which passes through the origin.

  The line $2x+by=6$ does not pass through the origin, so it cannot be the mirror line. For it to be invariant under the reflection, it must be perpendicular to the mirror line.

  The slope of $y=-2x$ is $-2$.

  Rewrite $2x+by=6$ as
  \[
  y=-\frac{2}{b}x+\frac{6}{b}
  \]
  so its slope is $-\dfrac{2}{b}$.

  Since perpendicular lines have slopes whose product is $-1$,
  \[
  (-2)\left(-\frac{2}{b}\right)=-1
  \]

  So
  \[
  \frac{4}{b}=-1
  \]
  giving
  \[
  b=-4
  \]

  Therefore $b=-4$.
\end{enumerate}
\end{tcolorbox}


\newpage

% ---- Question 4 ----
\questionitem
% [Source: _QBT___Solns__Invariant_Points_and_Lines, Q4]

The matrix $M$ is given by
\[
M = \begin{pmatrix} 1 & 0 \\[6pt] 0 & 6 \end{pmatrix}\begin{pmatrix} \frac{1}{2} & \frac{\sqrt{3}}{2} \\[6pt] \frac{\sqrt{3}}{2} & -\frac{1}{2} \end{pmatrix}
\] \\
\begin{enumerate}[label=\textbf{(\alph*)}]
  \item The matrix $M$ represents a sequence of two geometrical transformations in the $x$-$y$ plane. \\
  Give full details of the two transformations, stating clearly the order in which they are applied. \marks{4} \\
  \item Write $M^{-1}$ as the product of two matrices, neither of which is $I$. \marks{2} \\
  \item Find the equations of the invariant lines, through the origin, of the transformation represented by $M$. \marks{5} \\
\end{enumerate}

\vspace*{10pt}

\begin{tcolorbox}[colback=white, colframe=scarlet, title={\textbf{Solution}}, breakable, grow to left by=6mm, grow to right by=10mm]
\begin{enumerate}[label=\textbf{(\alph*)}]
\item Let
\[
D=\begin{pmatrix}1&0\\0&6\end{pmatrix},
\qquad
R=\begin{pmatrix}\frac12&\frac{\sqrt3}{2}\\[4pt]\frac{\sqrt3}{2}&-\frac12\end{pmatrix}
\]
so that
\[
M=DR
\]

First consider \(R\).  
A reflection in the line through the origin making angle \(\theta\) with the \(x\)-axis has matrix
\[
\begin{pmatrix}
\cos 2\theta & \sin 2\theta\\
\sin 2\theta & -\cos 2\theta
\end{pmatrix}
\]

Comparing with \(R\),
\[
\cos 2\theta=\frac12,
\qquad
\sin 2\theta=\frac{\sqrt3}{2}
\]
so
\[
2\theta=60^\circ
\quad\Rightarrow\quad
\theta=30^\circ
\]

Therefore \(R\) is a reflection in the line through the origin at \(30^\circ\) to the \(x\)-axis, whose equation is
\[
y=(\tan 30^\circ)x=\frac{x}{\sqrt3}
\]

Now consider \(D\). It maps
\[
\begin{pmatrix}x\\y\end{pmatrix}
\mapsto
\begin{pmatrix}x\\6y\end{pmatrix}
\]
so it is a stretch parallel to the \(y\)-axis with scale factor \(6\), with the \(x\)-axis invariant.

Since \(M=DR\), the matrix on the right acts first.

So the transformations are:

\[
\text{first reflect in } y=\frac{x}{\sqrt3},
\quad
\text{then stretch parallel to the } y\text{-axis with scale factor }6
\]

\item Using \(M=DR\),
\[
M^{-1}=(DR)^{-1}=R^{-1}D^{-1}
\]

A reflection is its own inverse, so
\[
R^{-1}=R
\]
and
\[
D^{-1}=\begin{pmatrix}1&0\\0&\frac16\end{pmatrix}
\]

Hence
\[
M^{-1}
=
\begin{pmatrix}
\frac12&\frac{\sqrt3}{2}\\[4pt]
\frac{\sqrt3}{2}&-\frac12
\end{pmatrix}
\begin{pmatrix}
1&0\\[4pt]
0&\frac16
\end{pmatrix}
\]

\item First find \(M\):
\[
M=
\begin{pmatrix}1&0\\0&6\end{pmatrix}
\begin{pmatrix}\frac12&\frac{\sqrt3}{2}\\[4pt]\frac{\sqrt3}{2}&-\frac12\end{pmatrix}
=
\begin{pmatrix}
\frac12&\frac{\sqrt3}{2}\\[4pt]
3\sqrt3&-3
\end{pmatrix}
\]

An invariant line through the origin must map to itself.

First check the vertical line \(x=0\). Its direction vector is \(\begin{pmatrix}0\\1\end{pmatrix}\), and
\[
M\begin{pmatrix}0\\1\end{pmatrix}
=
\begin{pmatrix}\frac{\sqrt3}{2}\\-3\end{pmatrix}
\]
which is not a multiple of \(\begin{pmatrix}0\\1\end{pmatrix}\), so \(x=0\) is not invariant.

So let an invariant line have equation
\[
y=mx
\]
with direction vector \(\begin{pmatrix}1\\m\end{pmatrix}\).

Its image is
\[
M\begin{pmatrix}1\\m\end{pmatrix}
=
\begin{pmatrix}
\frac12+\frac{\sqrt3}{2}m\\[4pt]
3\sqrt3-3m
\end{pmatrix}
\]

For the line to be invariant, this image must be parallel to \(\begin{pmatrix}1\\m\end{pmatrix}\), so
\[
\begin{pmatrix}
\frac12+\frac{\sqrt3}{2}m\\[4pt]
3\sqrt3-3m
\end{pmatrix}
=
k\begin{pmatrix}1\\m\end{pmatrix}
\]
for some constant \(k\).

Hence
\begin{align*}
\frac12+\frac{\sqrt3}{2}m &= k \\
3\sqrt3-3m &= km
\end{align*}

Eliminate \(k\):
\begin{align*}
3\sqrt3-3m &= m\left(\frac12+\frac{\sqrt3}{2}m\right) \\
6\sqrt3-6m &= m+\sqrt3 m^2 \\
\sqrt3 m^2+7m-6\sqrt3 &=0 \\
(\sqrt3 m-2)(m+3\sqrt3)&=0
\end{align*}

So
\[
m=\frac{2}{\sqrt3}
\qquad\text{or}\qquad
m=-3\sqrt3
\]

Therefore the invariant lines are
\[
y=\frac{2}{\sqrt3}x
\qquad\text{and}\qquad
y=-3\sqrt3 x
\]
\end{enumerate}
\end{tcolorbox}


\newpage

% ---- Question 5 ----
\questionitem
% [Source: _QBT___Solns__Invariant_Points_and_Lines, Q5]

A linear transformation of the plane is represented by the matrix
\[
\mathbf{M}=\begin{pmatrix}2 & -1\\ p & 0\end{pmatrix}
\]
where $p$ is a constant. \\
\begin{enumerate}[label=\textbf{(\alph*)}]
  \item Show that the line $x=0$ is not invariant. For a line through the origin with equation $y=mx$, show that it is invariant only if
  \[
  m^2-2m+p=0
  \]
  Hence find the set of values of $p$ for which the transformation has no invariant lines through the origin. \marks{5}\\
  \item Given that the transformation multiplies areas by $3$ and reverses orientation, find the invariant lines. \marks{3} \\
\end{enumerate}

\vspace*{10pt}

\begin{tcolorbox}[colback=white, colframe=scarlet, title={\textbf{Solution}}, breakable, grow to left by=6mm, grow to right by=10mm]
\begin{enumerate}[label=\textbf{(\alph*)}]
  \item A point on the line $x=0$ has the form
  \[
  \begin{pmatrix}0\\y\end{pmatrix}
  \]
  Applying the transformation,
  \[
  \mathbf{M}\begin{pmatrix}0\\y\end{pmatrix}
  =
  \begin{pmatrix}2&-1\\ p&0\end{pmatrix}
  \begin{pmatrix}0\\y\end{pmatrix}
  =
  \begin{pmatrix}-y\\0\end{pmatrix}
  \]
  This image has first coordinate $-y$, so it is not on the line $x=0$ unless $y=0$.

  Therefore the line $x=0$ is not invariant.

  Now consider a line through the origin with equation $y=mx$.

  A general point on this line is
  \[
  \begin{pmatrix}x\\mx\end{pmatrix}
  \]
  Its image is
  \[
  \mathbf{M}\begin{pmatrix}x\\mx\end{pmatrix}
  =
  \begin{pmatrix}2&-1\\ p&0\end{pmatrix}
  \begin{pmatrix}x\\mx\end{pmatrix}
  =
  \begin{pmatrix}2x-mx\\ px\end{pmatrix}
  =
  \begin{pmatrix}(2-m)x\\ px\end{pmatrix}
  \]
  For the line to be invariant, this image must also lie on $y=mx$.

  So its coordinates must satisfy
  \[
  px=m\bigl((2-m)x\bigr)
  \]
  Hence
  \[
  px=(2m-m^2)x
  \]
  so
  \[
  p=2m-m^2
  \]
  Rearranging,
  \[
  m^2-2m+p=0
  \]
  as required.

  The transformation has no invariant lines through the origin when this quadratic has no real solutions for $m$.

  Its discriminant is
  \[
  (-2)^2-4(1)(p)=4-4p
  \]
  For no real roots,
  \[
  4-4p<0
  \]
  \[
  1-p<0
  \]
  \[
  p>1
  \]

  Therefore the transformation has no invariant lines through the origin when $p>1$.

  \item The matrix of the transformation is
  \[
  \mathbf{M}=\begin{pmatrix}2&-1\\ p&0\end{pmatrix}
  \]
  Its determinant is
  \[
  \det(\mathbf{M})=2\cdot 0-(-1)(p)=p
  \]

  The transformation multiplies areas by $3$ and reverses orientation, so
  \[
  \det(\mathbf{M})=-3
  \]
  Hence
  \[
  p=-3
  \]

  Substitute this into the invariant-line condition from part (a):
  \[
  m^2-2m+p=0
  \]
  \[
  m^2-2m-3=0
  \]
  Factorising,
  \[
  (m-3)(m+1)=0
  \]
  So
  \[
  m=3 \quad \text{or} \quad m=-1
  \]

  Therefore the invariant lines are
  \[
  y=3x \quad \text{and} \quad y=-x
  \]
\end{enumerate}
\end{tcolorbox}


\newpage

% ---- Question 6 ----
\questionitem
% [Source: _QBT___Solns__Invariant_Points_and_Lines, Q6]

The $2 \times 2$ matrix $\mathbf{A}$ represents a transformation $\mathbf{T}$ which has the following properties. \\

\begin{itemize}
\item The image of the point $(1,-1)$ is the point $(2,0)$.
\item An object shape whose area is $4$ is transformed to an image shape whose area is $24$.
\item $\mathbf{T}$ has a line of invariant points. \\
\end{itemize}

\begin{enumerate}[label=(\roman*)]
  \item Find a possible matrix for $\mathbf{A}$. \marks{8} \\
\end{enumerate}

The transformation $\mathbf{S}$ is represented by the matrix $\mathbf{B}$ where
\[
\mathbf{B} = \begin{pmatrix}2 & 1 \\ 2 & 3\end{pmatrix}
\] \\
\begin{enumerate}[label=(\roman*), resume]
  \item Find the equation of the line of invariant points of $\mathbf{S}$. \marks{2} \\
  \item Show that any line of the form $y = 2x + c$ is an invariant line of $\mathbf{S}$. \marks{3} \\
\end{enumerate}

\vspace*{10pt}

\begin{tcolorbox}[colback=white, colframe=scarlet, title={\textbf{Solution}}, breakable, grow to left by=6mm, grow to right by=10mm]
\begin{enumerate}[label=\textbf{(\roman*)}]
\item Let
\[
\mathbf A=\begin{pmatrix}a&b\\ c&d\end{pmatrix}
\]

Using the fact that $(1,-1)$ maps to $(2,0)$,
\[
\mathbf A\begin{pmatrix}1\\-1\end{pmatrix}
=
\begin{pmatrix}2\\0\end{pmatrix}
\]

So
\begin{align*}
\begin{pmatrix}a&b\\ c&d\end{pmatrix}
\begin{pmatrix}1\\-1\end{pmatrix}
&=
\begin{pmatrix}a-b\\ c-d\end{pmatrix}
=
\begin{pmatrix}2\\0\end{pmatrix}
\end{align*}

Hence
\[
a-b=2,\qquad c-d=0
\]

So
\[
a=b+2,\qquad c=d
\]

Next, the area scale factor is
\[
\frac{24}{4}=6
\]

For a $2\times2$ matrix, the area scale factor is $|\det \mathbf A|$, so
\[
|\det \mathbf A|=6
\]

Now
\begin{align*}
\det \mathbf A&=ad-bc\\
&=(b+2)d-bd\\
&=2d
\end{align*}

Therefore
\[
|2d|=6
\]

So
\[
d=\pm3
\]

Now use the fact that $\mathbf T$ has a line of invariant points.

If a point on this line is $\begin{pmatrix}x\\y\end{pmatrix}$, then
\[
\mathbf A\begin{pmatrix}x\\y\end{pmatrix}=\begin{pmatrix}x\\y\end{pmatrix}
\]

Hence
\[
(\mathbf A-\mathbf I)\begin{pmatrix}x\\y\end{pmatrix}=\begin{pmatrix}0\\0\end{pmatrix}
\]

For there to be a whole line of invariant points, this equation must have non-zero solutions, so
\[
\det(\mathbf A-\mathbf I)=0
\]

Now
\[
\mathbf A-\mathbf I=
\begin{pmatrix}
a-1 & b\\
c & d-1
\end{pmatrix}
=
\begin{pmatrix}
b+1 & b\\
d & d-1
\end{pmatrix}
\]

So
\begin{align*}
\det(\mathbf A-\mathbf I)
&=(b+1)(d-1)-bd\\
&=bd-b+d-1-bd\\
&=d-b-1
\end{align*}

Setting this equal to $0$ gives
\[
d-b-1=0
\]

Hence
\[
b=d-1
\]

Then
\[
a=b+2=d+1,\qquad c=d
\]

If we choose $d=3$, then
\[
b=2,\qquad a=4,\qquad c=3
\]

So one possible matrix is
\[
\mathbf A=
\begin{pmatrix}
4&2\\
3&3
\end{pmatrix}
\]

A possible matrix is $\displaystyle \begin{pmatrix}4&2\\3&3\end{pmatrix}$.

\item Invariant points satisfy
\[
\mathbf B\begin{pmatrix}x\\y\end{pmatrix}=\begin{pmatrix}x\\y\end{pmatrix}
\]

So
\[
(\mathbf B-\mathbf I)\begin{pmatrix}x\\y\end{pmatrix}=\begin{pmatrix}0\\0\end{pmatrix}
\]

Now
\[
\mathbf B-\mathbf I=
\begin{pmatrix}
2&1\\
2&3
\end{pmatrix}
-
\begin{pmatrix}
1&0\\
0&1
\end{pmatrix}
=
\begin{pmatrix}
1&1\\
2&2
\end{pmatrix}
\]

Hence
\[
\begin{pmatrix}
1&1\\
2&2
\end{pmatrix}
\begin{pmatrix}x\\y\end{pmatrix}
=
\begin{pmatrix}0\\0\end{pmatrix}
\]

This gives
\[
x+y=0
\]

So the line of invariant points is
\[
y=-x
\]

\item Take any point on the line $y=2x+c$. Its coordinates are
\[
\begin{pmatrix}x\\2x+c\end{pmatrix}
\]

Its image under $\mathbf S$ is
\begin{align*}
\mathbf B\begin{pmatrix}x\\2x+c\end{pmatrix}
&=
\begin{pmatrix}
2&1\\
2&3
\end{pmatrix}
\begin{pmatrix}x\\2x+c\end{pmatrix}\\
&=
\begin{pmatrix}
2x+(2x+c)\\
2x+3(2x+c)
\end{pmatrix}\\
&=
\begin{pmatrix}
4x+c\\
8x+3c
\end{pmatrix}
\end{align*}

Let the image point be $(X,Y)$, where
\[
X=4x+c,\qquad Y=8x+3c
\]

Then
\begin{align*}
2X+c&=2(4x+c)+c\\
&=8x+2c+c\\
&=8x+3c\\
&=Y
\end{align*}

So the image satisfies
\[
Y=2X+c
\]

Therefore the image point also lies on the line $y=2x+c$.

Hence every line of the form $y=2x+c$ is an invariant line of $\mathbf S$.
\end{enumerate}
\end{tcolorbox}


\newpage

% ---- Question 7 ----
\questionitem
% [Source: _QBT___Solns__Invariant_Points_and_Lines, Q7]

The matrix
\[
B = \begin{pmatrix} 1 & -4 \\ 2 & -1 \end{pmatrix}
\]
represents the linear transformation $T$. \\
Any invariant line for $T$ must be parallel to an invariant line through the origin. \\
\begin{enumerate}[label=\textbf{(\alph*)}]
  \item A line through the origin has equation $y=mx$, where $m$ is real. Show that, for $m \ne \tfrac{1}{4}$, its image under $T$ has equation
  \[
  y = \frac{2-m}{1-4m}x
  \]
  State also the image of the line $y=\tfrac{1}{4}x$ and the image of the vertical line $x=0$. \marks{3}\\
  \item Hence prove that $T$ has no invariant lines. \marks{2} \\
\end{enumerate}

\vspace*{10pt}

\begin{tcolorbox}[colback=white, colframe=scarlet, title={\textbf{Solution}}, breakable, grow to left by=6mm, grow to right by=10mm]
\begin{enumerate}[label=\textbf{(\alph*)}]
  \item A general point on the line \(y=mx\) is
\[
\begin{pmatrix}x\\mx\end{pmatrix}
\]

Under the transformation \(T\) represented by
\[
B=\begin{pmatrix}1&-4\\2&-1\end{pmatrix}
\]
this point maps to
\[
\begin{pmatrix}X\\Y\end{pmatrix}
=
\begin{pmatrix}1&-4\\2&-1\end{pmatrix}
\begin{pmatrix}x\\mx\end{pmatrix}
=
\begin{pmatrix}x-4mx\\2x-mx\end{pmatrix}
=
\begin{pmatrix}(1-4m)x\\(2-m)x\end{pmatrix}
\]

If \(m\ne \tfrac14\), then \(1-4m\ne 0\), so
\[
\frac{Y}{X}
=
\frac{(2-m)x}{(1-4m)x}
=
\frac{2-m}{1-4m}
\]

Hence the image line has equation
\[
y=\frac{2-m}{1-4m}x
\]

Now take the special case \(m=\tfrac14\). A point on \(y=\tfrac14x\) is \(\begin{pmatrix}x\\x/4\end{pmatrix}\), and
\[
\begin{pmatrix}1&-4\\2&-1\end{pmatrix}
\begin{pmatrix}x\\x/4\end{pmatrix}
=
\begin{pmatrix}x-x\\2x-x/4\end{pmatrix}
=
\begin{pmatrix}0\\7x/4\end{pmatrix}
\]

So the image of \(y=\tfrac14x\) is
\[
x=0
\]

For the vertical line \(x=0\), a general point is \(\begin{pmatrix}0\\y\end{pmatrix}\). Its image is
\[
\begin{pmatrix}1&-4\\2&-1\end{pmatrix}
\begin{pmatrix}0\\y\end{pmatrix}
=
\begin{pmatrix}-4y\\-y\end{pmatrix}
\]

So
\[
Y=\frac14X
\]

Hence the image of \(x=0\) is
\[
y=\frac14x
\]

Therefore:
\[
y=mx \mapsto y=\frac{2-m}{1-4m}x \quad (m\ne \tfrac14), \qquad
y=\frac14x \mapsto x=0, \qquad
x=0 \mapsto y=\frac14x
\]

  \item For a line through the origin to be invariant, its image must have the same slope.

So for a non-vertical line \(y=mx\), we require
\[
m=\frac{2-m}{1-4m}
\]

Hence
\begin{align*}
m(1-4m)&=2-m\\
m-4m^2&=2-m\\
4m^2-2m+2&=0
\end{align*}

Its discriminant is
\[
(-2)^2-4(4)(2)=4-32=-28<0
\]

So there is no real value of \(m\), and therefore no invariant non-vertical line through the origin.

From part (a),
\[
y=\frac14x \mapsto x=0
\qquad\text{and}\qquad
x=0 \mapsto y=\frac14x
\]

So neither of these two lines is invariant either.

Therefore there are no invariant lines through the origin. Since any invariant line for \(T\) must be parallel to an invariant line through the origin, \(T\) has no invariant lines.
  \end{enumerate}
\end{tcolorbox}


\newpage

% ---- Question 8 ----
\questionitem
% [Source: _QBT___Solns__Invariant_Points_and_Lines, Q8]

The matrix of a linear transformation $T$ is
\[
M = \begin{pmatrix} p+2 & -5 \\[6pt] -2 & p-1 \end{pmatrix}
\] \\
\begin{enumerate}[label=\textbf{(\alph*)}]
  \item Find the values of $p$ for which $T$ is not invertible. \marks{2} \\
\end{enumerate}
It is now given that $p = 3$, so that
\[
M = \begin{pmatrix} 5 & -5 \\ -2 & 2 \end{pmatrix}
\] \\
\begin{enumerate}[label=\textbf{(\alph*)}, resume]
  \item Find the equations of the invariant lines, through the origin, of the transformation represented by $M$. \marks{6} \\
\end{enumerate}

\vspace*{10pt}

\begin{tcolorbox}[colback=white, colframe=scarlet, title={\textbf{Solution}}, breakable, grow to left by=6mm, grow to right by=10mm]
\begin{enumerate}[label=\textbf{(\alph*)}]
  \item A matrix is not invertible when its determinant is \(0\).

  \[
  \det(M)=
  \begin{vmatrix}
  p+2 & -5\\
  -2 & p-1
  \end{vmatrix}
  =(p+2)(p-1)-(-5)(-2)
  \]

  \[
  \det(M)=(p+2)(p-1)-10=p^2+p-2-10=p^2+p-12
  \]

  Factorising,

  \[
  p^2+p-12=(p+4)(p-3)
  \]

  So \(T\) is not invertible when

  \[
  (p+4)(p-3)=0
  \]

  Hence \(p=3\) or \(p=-4\).

  \item An invariant line through the origin must map onto itself.

  First check the vertical line \(x=0\). A general point on this line is \(\begin{pmatrix}0\\y\end{pmatrix}\), and

  \[
  M\begin{pmatrix}0\\y\end{pmatrix}
  =
  \begin{pmatrix}
  5 & -5\\
  -2 & 2
  \end{pmatrix}
  \begin{pmatrix}0\\y\end{pmatrix}
  =
  \begin{pmatrix}-5y\\2y\end{pmatrix}
  \]

  This does not in general have \(x=0\), so \(x=0\) is not invariant.

  Now take a non-vertical line through the origin with equation \(y=mx\).

  A general point on this line is \(\begin{pmatrix}x\\mx\end{pmatrix}\). Its image is

  \[
  M\begin{pmatrix}x\\mx\end{pmatrix}
  =
  \begin{pmatrix}
  5 & -5\\
  -2 & 2
  \end{pmatrix}
  \begin{pmatrix}x\\mx\end{pmatrix}
  =
  \begin{pmatrix}
  5x-5mx\\
  -2x+2mx
  \end{pmatrix}
  \]

  For the line to be invariant, this image must also lie on \(y=mx\). So its coordinates must satisfy

  \[
  -2x+2mx = m(5x-5mx)
  \]

  For a non-zero point, divide by \(x\):

  \[
  -2+2m = 5m-5m^2
  \]

  Rearranging,

  \[
  5m^2-3m-2=0
  \]

  Factorising,

  \[
  5m^2-3m-2=(5m+2)(m-1)
  \]

  So

  \[
  m=1 \qquad \text{or} \qquad m=-\frac25
  \]

  Therefore the invariant lines are

  \[
  y=x
  \]

  and

  \[
  y=-\frac25x
  \]

  Equivalently, the second line may be written as

  \[
  2x+5y=0
  \]

  So the invariant lines are \(y=x\) and \(2x+5y=0\).
\end{enumerate}
\end{tcolorbox}


\newpage

% ---- Question 9 ----
\questionitem
% [Source: _QBT___Solns__Invariant_Points_and_Lines, Q9]

\[
A = \begin{pmatrix} b & a - 3 \\[6pt] 2a - 1 & -2 \end{pmatrix}
\]
where $a$ and $b$ are non-zero constants. \\
Given that the matrix $A$ is self-inverse, \\
\begin{enumerate}[label=\textbf{(\alph*)}]
  \item determine the value of $b$ and the possible values of $a$. \marks{5} \\
\end{enumerate}
The matrix $A$ represents a linear transformation $M$ of the plane. \\
Using the smaller value of $a$ from part (a), \\
\begin{enumerate}[label=\textbf{(\alph*)}, resume]
  \item show that the invariant points of the linear transformation $M$ form a line, stating the equation of this line. \marks{3} \\
\end{enumerate}

\vspace*{10pt}

\begin{tcolorbox}[colback=white, colframe=scarlet, title={\textbf{Solution}}, breakable, grow to left by=6mm, grow to right by=10mm]
\begin{enumerate}[label=\textbf{(\alph*)}]
  \item Since \(A\) is self-inverse, \(A^{-1}=A\), so
\[
A^2=I
\]

Now
\[
A=\begin{pmatrix} b & a-3 \\ 2a-1 & -2 \end{pmatrix}
\]
so
\[
A^2=
\begin{pmatrix} b & a-3 \\ 2a-1 & -2 \end{pmatrix}
\begin{pmatrix} b & a-3 \\ 2a-1 & -2 \end{pmatrix}
=
\begin{pmatrix}
b^2+(a-3)(2a-1) & b(a-3)-2(a-3) \\
b(2a-1)-2(2a-1) & (2a-1)(a-3)+4
\end{pmatrix}
\]

Therefore
\[
\begin{pmatrix}
b^2+(a-3)(2a-1) & (a-3)(b-2) \\
(2a-1)(b-2) & (2a-1)(a-3)+4
\end{pmatrix}
=
\begin{pmatrix}
1 & 0 \\
0 & 1
\end{pmatrix}
\]

Equating entries gives
\[
(a-3)(b-2)=0
\]
\[
(2a-1)(b-2)=0
\]
and
\[
(2a-1)(a-3)+4=1
\]

Using the last equation,
\begin{align*}
(2a-1)(a-3)+4&=1\\
2a^2-7a+3+4&=1\\
2a^2-7a+6&=0\\
(2a-3)(a-2)&=0
\end{align*}

So
\[
a=\frac32 \quad \text{or} \quad a=2
\]

For both of these values, \(a\neq 3\) and \(a\neq \frac12\). Hence
\[
a-3\neq 0 \qquad \text{and} \qquad 2a-1\neq 0
\]

So from
\[
(a-3)(b-2)=0 \quad \text{and} \quad (2a-1)(b-2)=0
\]
it follows that
\[
b-2=0
\]
so
\[
b=2
\]

Hence
\[
b=2,\qquad a=\frac32 \text{ or } 2
\]

  \item Using the smaller value \(a=\frac32\), and \(b=2\),
\[
A=\begin{pmatrix} 2 & -\frac32 \\ 2 & -2 \end{pmatrix}
\]

Let \((x,y)\) be an invariant point. Then it is unchanged by the transformation, so
\[
\begin{pmatrix} 2 & -\frac32 \\ 2 & -2 \end{pmatrix}
\begin{pmatrix} x \\ y \end{pmatrix}
=
\begin{pmatrix} x \\ y \end{pmatrix}
\]

This gives the simultaneous equations
\begin{align*}
2x-\frac32 y&=x\\
2x-2y&=y
\end{align*}

So
\begin{align*}
x-\frac32 y&=0\\
2x-3y&=0
\end{align*}

These are the same equation, since the second is just twice the first. Therefore there are infinitely many invariant points, and they all lie on a straight line.

The line is
\[
2x-3y=0
\]

So the invariant points form the line \(2x-3y=0\), equivalently \(y=\frac{2}{3}x\).
\end{enumerate}
\end{tcolorbox}


\end{enumerate}
\end{document}
